% !TEX root = ../Monografia.tex
% ---------------------------------------------------------------------------------------------------------------------------------------------------------
% CAPÍTULO - CONTEXTO SOCIAL, AMBIENTAL E ÉTICO
% ---------------------------------------------------------------------------------------------------------------------------------------------------------

\chapter{Contexto Social, Ambiental e Ético}
\section{Impactos Socioeconômicos}

Eficiência computacional tem valor econômico mensurável. Em infraestrutura de servidores, a relação é direta: ciclos de CPU poupados são custos de nuvem reduzidos. Em dispositivos móveis, o mecanismo é diferente mas igualmente concreto: aplicações mais leves consomem menos bateria, reduzem a frequência de carregamentos e prolongam a vida útil dos dispositivos~\cite{cheng2020energy}. Em mercados como Brasil e Índia, onde o custo de um aparelho equivale a dois a três meses de salário mínimo, a longevidade do hardware é economicamente relevante para o usuário final.

Do ponto de vista social, aplicações mal otimizadas reforçam desigualdades de acesso. Um serviço de saúde ou de educação que funciona bem em um Pixel 8 mas gera \jank{} constante em um dispositivo de entrada cria uma experiência de segunda classe para usuários que já possuem menor capacidade de escolha. Otimizar desempenho, nesse contexto, também é uma questão de equidade no acesso a serviços digitais.

Ferramentas de software raramente são neutras em seus efeitos. Uma biblioteca de monitoramento de desempenho para Flutter parece, à primeira vista, um artefato técnico restrito ao ecossistema de desenvolvimento. Ainda assim, ela carrega decisões de projeto com implicações que vão além da eficiência computacional: para quem a ferramenta é acessível, quais dados ela coleta, como esses dados são tratados e em que medida ela apoia práticas de desenvolvimento mais conscientes do impacto ambiental e social das aplicações produzidas. Este capítulo situa o \texttt{collector\_flutter} nessas dimensões mais amplas.

\subsection*{Sustentabilidade Digital e o Custo Energético do Software}

A relação entre eficiência de software e consumo energético ganhou escala com a proliferação de dispositivos móveis. O relatório de sustentabilidade do Google~\cite{google2024sustainability} estima que melhorias na eficiência de renderização, gerenciamento de cache e políticas de rede podem reduzir o consumo energético de uma aplicação móvel típica em até 40\% sem degradação perceptível da experiência do usuário. O IEEE~\cite{ieee2023sustainability} situa o setor de TIC, incluindo dispositivos móveis, em aproximadamente 1,8\% das emissões globais de carbono, com projeção de crescimento sustentado ao longo da década.

Esses números têm uma implicação prática específica para o desenvolvimento de software: ineficiências locais se multiplicam pela escala. Um \textit{rebuild} desnecessário de widget que consome 2 ms de CPU extras, distribuído por 50 milhões de execuções diárias em uma base de usuários moderada, representa centenas de gigawatts-hora de energia desperdiçada por ano. Cheng e Wang~\cite{cheng2020energy} demonstraram que reduções de chamadas de rede redundantes e otimização de operações gráficas produzem economias energéticas mensuráveis mesmo em escala de dispositivo individual. Quando somado, esse efeito se torna ambientalmente relevante.

O \texttt{collector\_flutter} insere-se nessa perspectiva ao tornar visível, para o desenvolvedor, o custo computacional das escolhas de implementação. Ao identificar \textit{jank} recorrente, tendências de crescimento de memória ou volume excessivo de requisições de rede, a ferramenta ajuda a melhorar a experiência do usuário e a reduzir o consumo energético da aplicação resultante.

\section{Impactos Sociais e Inclusão Digital}

Desempenho não é uma preocupação uniforme ao longo da pirâmide socioeconômica. Aplicações que funcionam bem em dispositivos de alta capacidade, como Pixels, iPhones e Galaxy S, frequentemente degradam de forma severa em aparelhos de entrada que ainda representam a maioria absoluta do mercado em economias em desenvolvimento~\cite{google2024sustainability}. Essa assimetria tem consequências concretas: serviços de saúde, educação e comunicação que se tornam praticamente inutilizáveis em dispositivos de menor custo criam uma experiência de segunda classe para usuários que já possuem menor poder de escolha.

A otimização de desempenho em aplicações Flutter é, portanto, também uma questão de equidade de acesso. Uma aplicação que consome 40\% menos CPU e memória é mais inclusiva porque opera adequadamente em uma gama mais ampla de dispositivos. Ferramentas que facilitam a identificação e correção de ineficiências durante o desenvolvimento favorecem essa inclusão ao reduzir a distância entre uma aplicação funcionalmente correta e uma aplicação acessível para todos os usuários potenciais.

O \texttt{collector\_flutter} foi projetado para operar sem restrições de plataforma ou dispositivo. A coleta de métricas via \texttt{SchedulerBinding} e VM Service funciona igualmente em aparelhos de alta e baixa capacidade; as heurísticas de recomendação são calibradas para detectar comportamentos que afetam desempenho em dispositivos com restrições de hardware. Isso não torna a ferramenta um instrumento de inclusão digital por si só, mas alinha suas escolhas de projeto com esse valor.

\section{Impactos Ambientais e Eficiência Computacional}

O ciclo de vida de um dispositivo móvel tem dois determinantes principais: a durabilidade física do hardware e a compatibilidade do software disponível com suas especificações. Aplicações cada vez mais exigentes aceleram a obsolescência percebida dos dispositivos. Quando um smartphone de três anos não consegue executar as versões mais recentes dos aplicativos principais sem \textit{jank} severo, o incentivo para descartá-lo aumenta, mesmo que o hardware esteja fisicamente íntegro.

Resíduos eletrônicos constituem uma das correntes de lixo de crescimento mais rápido no mundo~\cite{unep2023}. Uma contribuição modesta, mas concreta, da engenharia de software para mitigar esse problema é o desenvolvimento de aplicações que continuem funcionando adequadamente em dispositivos mais antigos. Ferramentas de monitoramento que tornam visível o custo computacional do software e geram recomendações práticas para reduzi-lo apoiam esse objetivo.

\section{Aspectos Éticos no Monitoramento de Aplicações}

Sistemas de monitoramento de desempenho coletam informações sobre o comportamento de aplicações em execução. Em contextos de produção, onde o monitoramento ocorre sobre dispositivos de usuários reais, essa coleta levanta questões éticas legítimas: quais dados são coletados, por quanto tempo, por quem podem ser acessados e para quais finalidades podem ser utilizados.

O \texttt{collector\_flutter} foi projetado com o princípio de minimização de dados: ele coleta exclusivamente métricas de desempenho do sistema, como tempo de quadro, uso de memória, contagem de requisições e eventos definidos pelo desenvolvedor. Nenhum dado de identificação de usuário, conteúdo de comunicação ou informação pessoal é registrado. As métricas existem apenas em memória durante a sessão de uso da ferramenta e, na versão atual, não são transmitidas a servidores externos.

Essa delimitação é uma escolha ética e também técnica. Um sistema de monitoramento que coleta mais dados do que o necessário para seus objetivos de diagnóstico introduz riscos de privacidade desnecessários e aumenta o \textit{overhead} sem benefício correspondente para o desenvolvedor. O princípio de minimização de dados serve, portanto, tanto à ética quanto à eficiência.

A disponibilização do código-fonte como software aberto contribui adicionalmente para a transparência: qualquer desenvolvedor pode auditar o que é e o que não é coletado, verificar as heurísticas de recomendação e propor melhorias. Essa abertura é coerente com os padrões de reprodutibilidade científica recomendados por Wohlin \textit{et al.}~\cite{wohlin2012} para pesquisas experimentais em engenharia de software.

\section{Responsabilidade Tecnológica no Desenvolvimento de Software}

A noção de responsabilidade tecnológica, entendida como a corresponsabilidade de desenvolvedores e pesquisadores pelos impactos das soluções que produzem, é cada vez mais central na engenharia de software contemporânea. Na prática, isso significa considerar a correção funcional e o desempenho técnico de um sistema junto com suas implicações de acessibilidade, consumo de recursos, privacidade e longevidade.

Este trabalho é uma contribuição técnica: um pacote Dart com arquitetura bem definida e protocolo de avaliação experimental rigoroso. Ele também expressa uma escolha de projeto: o diagnóstico de desempenho deve ser acessível ao desenvolvedor que trabalha em uma equipe pequena, sem orçamento para ciclos de \textit{profiling} dedicados. A eficiência computacional deve ser monitorada continuamente, e ferramentas que democratizam esse diagnóstico ajudam a construir um ecossistema digital mais eficiente e mais equitativo.

\section{Síntese}

O monitoramento de desempenho em aplicações Flutter, enquanto problema técnico, é bem delimitado: coletar métricas, analisá-las e comunicar resultados ao desenvolvedor. Enquanto problema de engenharia com dimensões sociais e ambientais, é mais amplo: tornar esse diagnóstico acessível reduz o consumo energético das aplicações produzidas, amplia sua compatibilidade com dispositivos de menor capacidade e incentiva práticas de desenvolvimento mais conscientes de seu impacto. O \texttt{collector\_flutter} é uma resposta ao primeiro problema e busca apoiar, de forma instrumental, o segundo.
